%-------------------------
% Resume in Latex -- TCS Prime interview variant
% Based on Jake's Resume Template (https://github.com/jakegut/resume / sourabhbajaj)
%
% Differences from resume_26_aug.tex, and why:
%  - Class X and XII rows added to Education. TCS's eligibility frame is
%    10th/12th/UG and panels ask for those numbers directly. Marks here must
%    match the iBegin/NextStep portal entries exactly or verification fails.
%  - No Objective section. It fits only at a 0.15in bottom margin, which is
%    inside the unprintable region of many printers, and this variant exists
%    to be printed. Dropping it leaves a 0.45in bottom margin. The content
%    lives in the spoken "tell me about yourself" script instead. Re-adding it
%    costs one RoPE-DiT bullet.
%  - Core CS line added to skills. TCS TR is OS/DBMS/CN/OOP heavy; naming them
%    steers the panel onto ground that can be revised the night before.
%  - Skills cut from ~60 items to ~30. Every item removed had zero occurrences
%    across all nine repos: LangChain, LangGraph, CrewAI, pgvector, llama.cpp,
%    Unsloth, TRL, NCNN, Flask, Spark, CUDA, SageMaker, Redshift, EMR, ECS/ECR,
%    CloudWatch, BigQuery, Dataproc, Dataflow, Pub/Sub, Vertex AI, vLLM.
%    A TCS panel picks items off this list at random. Do not re-add one
%    without code to open.
%  - Bullets shortened to roughly 25 words. Long bullets cannot be spoken
%    aloud, and every bullet here has to survive being narrated in the room.
%
% ATS notes:
%  - Single column, standard section headings, no text boxes, no images.
%    Selectable text with \pdfgentounicode=1 so parsers read it cleanly.
%  - Link text uses ulem's \uline, not \underline: \underline typesets in math
%    mode, which ignores \small and extracts worse. Do not switch it back.
%  - Every claim below is backed by code that exists in a repo. Do not re-add
%    a bullet you cannot open a file for.
%------------------------

\documentclass[letterpaper,11pt]{article}

\usepackage{latexsym}
\usepackage[empty]{fullpage}
\usepackage{titlesec}
\usepackage{marvosym}
\usepackage[usenames,dvipsnames]{color}
\usepackage{verbatim}
\usepackage{enumitem}
\usepackage[hidelinks]{hyperref}
\usepackage{fancyhdr}
\usepackage[english]{babel}
\usepackage[normalem]{ulem}
\usepackage{tabularx}
\input{glyphtounicode}

\pagestyle{fancy}
\fancyhf{}
\fancyfoot{}
\renewcommand{\headrulewidth}{0pt}
\renewcommand{\footrulewidth}{0pt}

% Adjust margins
\addtolength{\oddsidemargin}{-0.7in}
\addtolength{\evensidemargin}{-0.7in}
\addtolength{\textwidth}{1.4in}
\addtolength{\topmargin}{-0.7in}
\addtolength{\textheight}{1.55in}

\urlstyle{same}

\raggedbottom
\raggedright
\setlength{\tabcolsep}{0in}

% Sections formatting
\titleformat{\section}{
  \vspace{-8pt}\scshape\raggedright\large
}{}{0em}{}[\color{black}\titlerule \vspace{-8pt}]

% Ensure that generated pdf is machine readable / ATS parsable
\pdfgentounicode=1

\newcommand{\graddate}{2023 -- 2027}

\newcommand{\introle}{Machine Learning Intern}
\newcommand{\intdates}{2026 -- Present}

%-------------------------
% Custom commands
\newcommand{\resumeItem}[1]{
  \item\small{#1}
}

\newcommand{\resumeSubheading}[4]{
  \vspace{-2pt}\item
    \begin{tabular*}{0.97\textwidth}[t]{l@{\extracolsep{\fill}}r}
      \textbf{#1} & #2 \\
      \textit{\small#3} & \textit{\small #4} \\
    \end{tabular*}\vspace{-8pt}
}

% Single-line education row, for the school rows where a second line would
% waste a third of the Education section on two numbers.
\newcommand{\resumeEduLine}[4]{
  \vspace{1pt}\item
    \begin{tabular*}{0.97\textwidth}[t]{l@{\extracolsep{\fill}}r}
      \textbf{#1}, \textit{\small#2} & \textit{\small#3} \hspace{6pt} #4 \\
    \end{tabular*}\vspace{-8pt}
}

\newcommand{\resumeProjectHeading}[2]{
    \item
    \begin{tabular*}{0.97\textwidth}{l@{\extracolsep{\fill}}r}
      \small#1 & \small#2 \\
    \end{tabular*}\vspace{-5pt}
}

\newcommand{\resumeSubItem}[1]{\resumeItem{#1}\vspace{-4pt}}

\renewcommand\labelitemii{$\vcenter{\hbox{\tiny$\bullet$}}$}

\newcommand{\resumeSubHeadingListStart}{\begin{itemize}[leftmargin=0.15in, label={}, itemsep=0pt]}
\newcommand{\resumeSubHeadingListEnd}{\end{itemize}\vspace{-11pt}}
\newcommand{\resumeItemListStart}{\begin{itemize}[itemsep=2pt, parsep=0pt, topsep=0pt]}
\newcommand{\resumeItemListEnd}{\end{itemize}\vspace{-6pt}}

%-------------------------------------------
%%%%%%  RESUME STARTS HERE  %%%%%%%%%%%%%%%%%%%%%%%%%%%%%%%%%%%%%%%%

\begin{document}

%----------HEADING----------
\begin{center}
    \textbf{\Huge \scshape Divyansh Shukla} \\ \vspace{2pt}
    \scriptsize
    \href{mailto:divyanshmohanshukla@hotmail.com}{\uline{divyanshmohanshukla@hotmail.com}} $|$
    +91 7827910144 $|$
    \href{https://github.com/divyanshuklai}{\uline{github.com/divyanshuklai}} $|$
    \href{https://www.linkedin.com/in/divyansh-shukla-ai/}{\uline{linkedin.com/in/divyansh-shukla-ai}} $|$
    \href{https://divyanshuklai.github.io}{\uline{divyanshuklai.github.io}}
\end{center}

%-----------EDUCATION-----------
\section{Education}
  \resumeSubHeadingListStart
    \resumeSubheading
      {Vellore Institute of Technology, Andhra Pradesh}{\graddate}
      {B.Tech, Computer Science and Engineering (Artificial Intelligence \& Machine Learning)}{CGPA: 9.03 / 10}
    \resumeEduLine{DAV Public School, Sahibabad}{Class XII, CBSE}{2022}{91\%}
    \resumeEduLine{DAV Public School, Sahibabad}{Class X, CBSE}{2020}{95.6\%}
    \vspace{3pt}
  \resumeSubHeadingListEnd

%-----------SKILLS-----------
\section{Technical Skills}
 \begin{itemize}[leftmargin=0.15in, label={}]
    \small{\item{
     \textbf{Languages}{: Python, SQL, C++, JavaScript} \\
     \textbf{Core CS}{: Data Structures and Algorithms, DBMS, Operating Systems, Computer Networks, Object-Oriented Programming, SDLC and Agile} \\
     \textbf{Web \& Backend}{: FastAPI, Django, REST APIs, PostgreSQL, HTML, CSS} \\
     \textbf{Tools \& Infrastructure}{: Git, Linux, Docker, Kubernetes (basics), AWS (EC2, S3, Lambda, IAM), Modal, Weights \& Biases, Hydra} \\
     \textbf{ML Frameworks}{: PyTorch, PyTorch Lightning, Hugging Face Transformers, Ultralytics (YOLO), Stable-Baselines3, NumPy, pandas, OpenCV} \\
     \textbf{Deep Learning}{: object detection, flow and diffusion models, transformers, reinforcement learning (PPO)} \\
     \textbf{Generative AI}{: Google Gen AI SDK (Gemini), hand-written agent loops (planning, memory, function calling), RAG with ChromaDB}
    }}
 \end{itemize}

%-----------EXPERIENCE-----------
\section{Experience}
  \resumeSubHeadingListStart
    \resumeSubheading
      {Power Pulse India}{\intdates}
      {\introle}{Ghaziabad, India}
      \resumeItemListStart
        \resumeItem{Trained a YOLOv8n thermal object detector on FLIR ADAS v2 across 6 road-safety classes. Cut missed pedestrians by 26\% and raised mAP50-95 from 0.360 to 0.420, giving up 3 points of precision for the recall a safety case needs.}
        \resumeItem{Audited per-class crops and COCO size buckets to find the accuracy ceiling. 82\% of pedestrian labels are under 32x32 px, and 37\% of true trucks drew no prediction at all.}
        \resumeItem{Ran a four-experiment resolution and architecture sweep. Dropped a P2 detection head that scored below plain 640px at 1.51x the compute, and chose 800px over 960px.}
      \resumeItemListEnd
  \resumeSubHeadingListEnd

%-----------PROJECTS-----------
\section{Projects}
    \resumeSubHeadingListStart

      \resumeProjectHeading
          {\textbf{RoPE-DiT: Diffusion Transformer with 2D Rotary Embeddings}}{\href{https://github.com/divyanshuklai/Custom2DRoPEDiT}{\uline{github.com/divyanshuklai/Custom2DRoPEDiT}}}
          \resumeItemListStart
            \resumeItem{Wrote a 5M-parameter Diffusion Transformer from scratch in PyTorch: patchifier, timestep and class-label embedders, AdaLN-Zero conditioning, and a custom multi-head attention module.}
            \resumeItem{Built a 2D rotary positional embedding that rotates queries and keys by each patch's row and column, so attention sees relative 2D distance instead of a flat index.}
            \resumeItem{Trained on CIFAR-10 with a rectified flow objective. Sampled with a hand-written Euler solver and classifier-free guidance.}
          \resumeItemListEnd

      \resumeProjectHeading
          {\textbf{Fashion Search: Conversational Product Discovery}}{\href{https://github.com/divyanshuklai/fashion_search}{\uline{github.com/divyanshuklai/fashion\_search}}}
          \resumeItemListStart
            \resumeItem{Built a FastAPI chat app that answers natural-language product queries across 10 Indian e-commerce stores, using Tavily search, Gemini 2.5 Flash extraction and BeautifulSoup scraping.}
            \resumeItem{Wrote the agent loop by hand on the Google Gen AI SDK, with no agent framework: a planner step decides whether to search or ask a clarifying question, conversation history acts as memory, and Gemini function calls are dispatched manually.}
            \resumeItem{Wrote a priority-based price and image extractor that tries JSON-LD product schema first, then Open Graph tags, then keyword-matched image tags.}
          \resumeItemListEnd

      \resumeProjectHeading
          {\textbf{Zero-Shot ICE RMA: Quadruped Locomotion RL}}{\href{https://github.com/divyanshuklai/zero-shot-ice-rma}{\uline{github.com/divyanshuklai/zero-shot-ice-rma}}}
          \resumeItemListStart
            \resumeItem{Built a MuJoCo environment for the Unitree Go1 with 30-dim proprioceptive observations, 12-dim joint-position actions and fractal terrain, then debugged a PPO baseline that ran 42M steps without learning a gait and traced the exploding value loss to unnormalized observations.}
          \resumeItemListEnd

    \resumeSubHeadingListEnd

%-----------CERTIFICATIONS-----------
\section{Certifications \& Achievements}
 \begin{itemize}[leftmargin=0.15in, label={}]
    \small{\item{
     \textbf{NPTEL}{: Deep Learning for Computer Vision (Topper, top 1\% of 463, IIT-H); Applied Accelerated AI (Topper, top 2\% of 591, IIT-G)} \\
     \textbf{Coursera}{: Deep Learning Specialization} \hspace{3pt}$|$\hspace{3pt} \textbf{AWS}{: AWS Academy Cloud Architecting, AWS Cloud Practitioner (badge)} \\
     \textbf{Achievement}{: Won the intra-campus round of the Smart India Hackathon at VIT-AP, 2024}
    }}
 \end{itemize}

\end{document}
